%!TEX root = forallxyyc.tex
\part{元理论}
\label{ch.normalform}
\addtocontents{toc}{\protect\mbox{}\protect\hrulefill\par}

\chapter{范式}\label{c:NormalForms}

\section{析取范式}\label{s:DNFDefined}

有时，考虑形式特别简单的语句是很有用的。例如，我们可以考虑那些 $\enot$ 只出现在原子语句之前的语句，或者那些只用 $\eand$ 将原子语句和带否定的原子语句组合起来的语句。一种相对通用但仍保持简单的形式是：一个语句是若干原子语句或带否定的原子语句的合取式的析取。当构造这样一个语句时，我们从原子语句开始，然后（可能）附加否定，接着（可能）用 $\eand$ 组合，最后（可能）用 $\eor$ 组合。

我们说一个语句是 \define{析取范式（disjunctive normal form）} \emph{\ifeff} 它满足以下所有条件：

\begin{compactlist}
		\item[(\textsc{dnf1})] 除了否定、合取和析取外，语句中不出现其他联结词；
		\item[(\textsc{dnf2})] 每个否定的出现都具有最小辖域（即，任何“$\enot$”都紧跟着一个原子语句）；
		\item[(\textsc{dnf3})] 任何析取的辖域内都不出现合取。
	\end{compactlist}

\newglossaryentry{disjunctive normal form}{
  name = 析取范式,
  text = 析取范式,
  description = {一种语句，它是若干原子语句或带否定的原子语句的合取式的析取}
}
所以，以下是一些属于析取范式的语句：
\begin{align*}
  & A \\
  & (A \eand \lnot B \eand C)\\
  & (A \eand B) \eor (A \eand \enot B)\\
  & (A \eand B) \eor (A \eand  B \eand C \eand \enot D \eand \enot E)\\
  & A \eor (C \eand \enot P_{234} \eand P_{233} \eand Q) \eor \enot B
\end{align*}
注意，我们这里暂时打破了本书的一项原则，\emph{临时}允许自己采用宽松的括号约定，允许合取式和析取式具有任意长度。这些约定使得判断一个语句是否属于析取范式变得更加容易。后续我们将继续采用这些宽松约定，不再另行说明。

为了进一步阐明析取范式的概念，我们将引入更多记号。我们用“$\pm\metav{A}$”表示 $\metav{A}$ 是一个原子语句，它前面可能带有一个否定符号，也可能没有。那么，一个属于析取范式的语句具有以下形状：
	$$(\pm \metav{A}_1 \land \ldots \land \pm \metav{A}_i) \lor (\pm \metav{A}_{i+1} \land \ldots \land \pm\metav{A}_j) \lor \ldots \lor (\pm\metav{A}_{m+1} \land \ldots \land \pm \metav{A}_n)$$
我们现在知道了什么是属于析取范式的语句。我们旨在证明的结果是：
	\factoidbox{\label{thm:dnf}\textbf{析取范式定理。} 对任意语句，都存在一个与之逻辑等价的、属于析取范式的语句。
	}
\section{通过真值表证明析取范式定理}
\label{s:DNFTruthTable}

我们对析取范式定理的第一个证明使用了真值表。我们将首先举例说明如何找出等价的析取范式语句，然后再将这一示例转化为严格的证明。

假设我们有一个语句 $\metav{S}$，它包含三个原子语句“$A$”、“$B$” 和 “$C$”。首先要做的是为 $\metav{S}$ 画出一张完整的真值表。结果可能如下：

\begin{center}
\begin{tabular}{c c c | c}
$A$ & $B$ & $C$ & $\metav{S}$\\
\hline
 T & T & T & T \\
 T & T & F & F \\
 T & F & T & T \\
 T & F & F & F \\
 F & T & T & F \\
 F & T & F & F \\
 F & F & T & T \\
 F & F & F & T
\end{tabular}
\end{center}

%Now, consider a sentence, whose only connectives are negations and conjunctions, where no connective occurs within the scope of any negation, e.g.:
%	$$A \eand \enot B \eand C$$
%This sentence is true when, and only when, “$A$” is true, “$B$” is false and “$C$” is true. Similarly, the sentence:
%	$$\enot A \eand \enot B \eand C$$
%this is true when, and only when, “$A$” is false, “$B$” is false and “$C$” is true.
%
%A disjunction is true when, and only when, at least one of the disjuncts is true. So if we write down a disjunction of sentences of the above form, perhaps
%	$$(A \eand \enot B \eand C) \eor (\enot A \eand \enot B \eand C)$$
%then it will be true on exactly \emph{two} lines of the truth table which describes all possible valuations of “$A$”, “$B$” and “$C$”.
%
恰巧，$\metav{S}$ 在其真值表的四行上为真，即第 1、3、7 和 8 行。对应于每一行，我们将写下四个只包含否定和合取作为联结词的语句，并且其中每个否定的辖域都保持最小：

\begin{compactlist}
		\item  “$A \eand B \eand C$”\hfill 该式在第 1 行为真（且仅当该行）
		\item “$A \eand \enot B \eand C$” \hfill 该式在第 3 行为真（且仅当该行）
		\item “$\enot A \eand \enot B \eand C$” \hfill 该式在第 7 行为真（且仅当该行）
		\item “$\enot A \eand \enot B \eand \enot C$” \hfill 该式在第 8 行为真（且仅当该行）
	\end{compactlist}

现在我们用 $\eor$ 将所有合取式组合起来：
$$(A \eand B \eand C) \eor (A \eand \enot B \eand C) \eor (\enot A \eand \enot B \eand C) \eor (\enot A \eand \enot B \eand \enot C)$$\label{longDNF}
这样就得到了一个析取范式语句，它恰好在那些至少有一个析取支为真的行上为真，即在（且仅在）第 1、3、7 和 8 行为真。因此，这个语句与 $\metav{S}$ 具有完全相同的真值表。于是我们得到了一个与 $\metav{S}$ 逻辑等价的析取范式语句，这正是我们想要的！

我们刚才采用的策略并不依赖于 $\metav{S}$ 的具体细节；它是完全通用的。因此，我们可以利用它来给出析取范式定理的一个简单证明。

任选一个语句 $\metav{S}$，并令 $\metav{A}_1, \ldots, \metav{A}_n$ 为出现在 $\metav{S}$ 中的所有原子语句。为得到一个逻辑等价于 $\metav{S}$ 的析取范式语句，我们需要考察 $\metav{S}$ 的真值表。有两种情况需要考虑：

\begin{numberlist}
		\item \emph{$\metav{S}$ 在其真值表的每一行上均为假。} 那么 $\metav{S}$ 是一个矛盾式。在这种情况下，矛盾式 $(\metav{A}_1 \eand \enot \metav{A}_1)$ 是一个析取范式语句，且逻辑等价于 $\metav{S}$。
		\item \emph{$\metav{S}$ 在其真值表的至少一行上为真。}
		对于真值表的每一行 $i$，令 $\metav{B}_i$ 为一个形如
		$$(\pm\metav{A}_1 \land \ldots \land \pm\metav{A}_n)$$
		的合取式，其中决定是否在每一个原子语句前添加否定的规则如下：
			\begin{itemize}
				\item $\metav{A}_m$ 是 $\metav{B}_i$ 的一个合取支 \emph{\ifeff} $\metav{A}_m$ 在第 $i$ 行为真。
				\item $\enot\metav{A}_m$ 是 $\metav{B}_i$ 的一个合取支 \emph{\ifeff} $\metav{A}_m$ 在第 $i$ 行为假。
			\end{itemize}
		根据这些规则，$\metav{B_i}$ 在且仅在真值表考虑 $\metav{A}_1$, \ldots, $\metav{A}_n$ 所有可能赋值（即 $\metav{S}$ 的真值表）的第 $i$ 行上为真。

		接下来，令 $i_1$, $i_2$, \dots, $i_m$ 为真值表中 $\metav{S}$ \emph{为真} 的那些行的编号。现在令 $\metav{D}$ 为如下语句：
		$$\metav{B}_{i_1} \eor \metav{B}_{i_2} \eor \ldots \eor \metav{B}_{i_m}$$
		既然 $\metav{S}$ 在其真值表的至少一行上为真，那么 $\metav{D}$ 就是良定义的；在极限情况下，即 $\metav{S}$ 在其真值表上恰好只在一行上为真时，$\metav{D}$ 就直接是某个 $\metav{B}_{i_1}$。

		根据构造，$\metav{D}$ 是析取范式语句。此外，根据构造，对于真值表的每一行 $i$：$\metav{S}$ 在该行上为真 \emph{\ifeff} $\metav{D}$ 的某个析取支（即 $\metav{B_i}$）在且仅在该行上为真。因此，$\metav{S}$ 和 $\metav{D}$ 具有相同的真值表，从而逻辑等价。
	\end{numberlist}

	这两种情况已穷尽所有可能，无论如何，我们都得到了一个逻辑等价于 $\metav{S}$ 的析取范式语句。

于是我们证明了析取范式定理。不过，在我们进一步讨论之前，应当立即指出，我们因此回到了对（真值函项逻辑）语句的严格定义，根据该定义，可以假定任一合取式恰有两个合取支，任一个析取式恰有两个析取支。

\section{合取范式}
\label{s:CNF}

本章至此我们讨论的是\emph{析取}范式。存在所谓的\emph{合取范式}，可能并不令人意外。

合取范式 的定义与析取范式的定义完全类似。确切地说，一个语句是合取范式语句 \emph{\ifeff} 它满足以下所有条件：

\begin{compactlist}
		\item[(\textsc{cnf1})] 除了否定、合取和析取之外，该语句中不出现任何其他联结词；
		\item[(\textsc{cnf2})] 每个否定的出现都具有最小辖域；
		\item[(\textsc{cnf3})] 任何合取的辖域内都不出现析取。
	\end{compactlist}

\newglossaryentry{conjunctive normal form}{
  name = 合取范式,
  text = 合取范式,
  description = {一种语句，它是若干原子语句或带否定的原子语句的析取式的合取}
}
一般而言，一个合取范式语句形如
	$$(\pm \metav{A}_1 \lor \ldots \lor \pm \metav{A}_i) \land (\pm \metav{A}_{i+1} \lor \ldots \lor \pm\metav{A}_j) \land \ldots \land (\pm\metav{A}_{m+1} \lor\ldots \lor \pm \metav{A}_n)$$
其中每个 $\metav{A}_k$ 都是一个原子语句。

我们现在可以证明另一个范式定理：
	\factoidbox{\label{thm:cnf}\textbf{合取范式定理。} 对于任一语句，都存在一个与其逻辑等价的合取范式语句。}

	给定一个真值函项逻辑语句 $\metav{S}$，我们首先写下 $\metav{S}$ 的完整真值表。

	如果 $\metav{S}$ 在真值表的\emph{每一行}上都为真，那么 $\metav{S}$ 与 $(\metav{A}_1 \eor \enot \metav{A}_1)$ 逻辑等价。

	如果 $\metav{S}$ 在真值表的\emph{至少一行}上为假，那么，对于真值表中 $\metav{S}$ 为假的每一行，写下在该行上（且仅在该行上）为\emph{假}的析取式 $(\pm\metav{A}_1 \eor \ldots \eor \pm\metav{A}_n)$。令 $\metav{C}$ 为所有这些析取式的合取；根据构造，$\metav{C}$ 是合取范式语句，并且 $\metav{S}$ 与 $\metav{C}$ 逻辑等价。

\practiceproblems
\problempart
\label{pr.DNF}
考虑以下语句：

\begin{compactlist}
		\item $(A \eif \enot B)$
		\item $\enot (A \eiff B)$
		\item $(\enot A \eor \enot (A \eand B))$
		\item $(\enot (A \eif B ) \eand (A \eif C))$
		\item $(\enot (A \eor B) \eiff ((\enot C \eand \enot A) \eif \enot B))$
		\item $((\enot (A \eand \enot B) \eif C) \eand \enot (A \eand D))$
	\end{compactlist}

        对每一个语句，分别找出一个逻辑等价的析取范式语句和一个逻辑等价的合取范式语句。

\chapter{功能完全性}\label{c:FunctionalCompleteness}

在我们的联结词中，$\enot$ 仅作用于一个句子，其他联结词都恰好组合两个句子。我们也可以引入 $n$ 元联结词的概念。例如，我们可以考虑一个三元联结词 “$\heartsuit$”，并规定它具有以下特征真值表：

\begin{center}
\begin{tabular}{c c c | c}
$A$ & $B$ & $C$ & $\heartsuit(A,B,C)$\\
\hline
 T & T & T & F \\
 T & T & F & T \\
 T & F & T & T \\
 T & F & F & F \\
 F & T & T & F \\
 F & T & F & T \\
 F & F & T & F \\
 F & F & F & F
\end{tabular}
\end{center}

可能这个新联结词不对应任何自然的英语表达式（至少不像 “$\eand$” 对应“and”那样）。但问题随之产生：如果我们想使用具有这个特征真值表的联结词，是否必须在真值函项逻辑中添加一个\emph{新}联结词？还是说我们可以用\emph{已有}的联结词来达到目的（就像我们可以用已有联结词表达“既不……也不……”那样）？

让我们更精确地表述这个问题。我们说一组联结词是\define{联合功能完备的} \emph{\ifeff}，对于任何可能的真值表，都存在一个只包含这些联结词的句子具有该真值表。

\newglossaryentry{functionally complete}{
  name = {功能完备性},
  text = {功能完备的},
  description = {联结词集合的一个性质，该性质成立 \ifeff{} 每个可能的真值表都是某个仅包含那些联结词的句子的真值表}}

一般要点是，当我们掌握了一些联合功能完备的联结词，就没有哪个特征真值表是我们无法把握的。事实上，我们很幸运。
	\factoidbox{\label{thm:ExpressiveAdequacy}\textbf{功能完备性定理。}
真值函项逻辑的联结词是联合功能完备的。事实上，以下几对联结词都是联合功能完备的：

\begin{compactlist}
\item\label{expressive:eor} “$\enot$” 和 “$\eor$”
\item\label{expressive:eand} “$\enot$” 和 “$\eand$”
\item\label{expressive:eif} “$\enot$” 和 “$\eif$”
\end{compactlist}
}

给定任何真值表，我们可以使用\cref{c:NormalForms}中通过真值表证明析取范式定理（或合取范式定理）的方法，写下一个具有相同真值表的模式。例如，运用证明析取范式定理的真值表方法，我们发现以下模式具有与上面$\heartsuit(A,B,C)$相同的特征真值表：
		$$(A \eand B \eand \enot C) \eor (A \eand \enot B \eand C) \eor (\enot A \eand B \eand \enot C)$$
由此可见，真值函项逻辑的联结词是联合功能完备的。我们现在证明每一个辅助结果。

\emph{辅助结果1：“$\enot$”和“$\eor$”的功能完备性。}注意，使用证明析取范式定理的真值表方法生成的模式将只包含联结词“$\enot$”、“$\eand$”和“$\eor$”。因此，只需证明存在一个只包含“$\enot$”和“$\eor$”的等价模式即可。为达此目的，我们只需注意到：
		\begin{align*}
		(\metav{A} \eand \metav{B}) & \Leftrightarrow \enot(\enot \metav{A} \eor\enot \metav{B})
		\end{align*}
		是逻辑等价的。\footnote{这里我们采用一个约定：称$\metav{A}$和$\metav{B}$等价可缩写为$\metav{A} \Leftrightarrow \metav{B}$。注意“$\Leftrightarrow$” \emph{不是}像“$\eiff$”那样的联结词，而是像“$\entails$”一样的元语言符号；比较\cref{entails-v-iff}中对“$\entails$”和“$\eiff$”的讨论。}

\emph{辅助结果2：“$\enot$”和“$\eand$”的功能完备性。}与辅助结果1完全相同，利用以下事实：
		\begin{align*}
		(\metav{A} \eor \metav{B}) & \Leftrightarrow \enot(\enot \metav{A} \eand\enot \metav{B})
		\end{align*}
		是逻辑等价的。

\emph{辅助结果3：“$\enot$”和“$\eif$”的功能完备性。}与辅助结果1完全相同，但改用以下等价式：
		\begin{align*}
		(\metav{A} \eor \metav{B}) &\Leftrightarrow (\enot \metav{A} \eif \metav{B})\\
		(\metav{A} \eand \metav{B}) &\Leftrightarrow \enot(\metav{A} \eif \enot\metav{B})
		\end{align*}
或者，我们可以直接依赖另外两个辅助结果之一，并且（反复地）只调用这两个等价式中的一个。

总之，我们\emph{从不必须}在真值函项逻辑中添加新的联结词。事实上，我们已有的联结词中已经存在冗余：如果我们当初真的力求精简，本可以只用两个联结词就够了。

\ifHTMLtarget
\section{单独功能完备的联结词}
\else
\section[单独功能完备的联结词][功能完备的联结词]{单独功能完备的联结词}
\fi

事实上，一些二元联结词是\emph{单独}功能完备的。这些联结词通常不包含在真值函项逻辑中，因为它们使用起来相当繁琐。但它们的存在表明，倘若我们愿意，本可以定义一种功能完备的真值函项语言，它仅包含一个初始联结词。

我们将考虑的第一个此类联结词是 “$\uparrow$”，它具有以下特征真值表。

\begin{center}
\begin{tabular}{c c | c}
$\metav{A}$ & $\metav{B}$ & $\metav{A} \mathrel{\uparrow} \metav{B}$\\
\hline
 T & T & F \\
 T & F & T \\
 F & T & T  \\
 F & F & T
\end{tabular}
\end{center}

这通常被称为“谢弗竖线”，得名于亨利·谢弗，他曾用它来展示如何减少 Russell 和 Alfred North Whitehead（阿尔弗雷德·诺思·怀特海）《数学原理》中的逻辑联结词数量。\footnote{Sheffer, “”A set of five independent postulates for Boolean algebras, with application to logical constants'', \textit{Transactions of the American Mathematical Society}~14 (1913), pp. 481--488}（实际上，查尔斯·桑德斯·皮尔士比谢弗早了大约30年就预见到了这一点，但他从未发表其结果；而波兰哲学家爱德华·斯塔姆则比谢弗早两年发表了同样的结果。）\footnote{参见皮尔士约1880年的文章“A Boolian algebra with one constant”，收录于Peirce's \textit{Collected Papers}, vol. 4, pp. 264--5；以及爱德华·斯塔姆，“”\foreignlanguage{german}{Beitrag zur Algebra der Logik}'', \foreignlanguage{german}{\textit{Monatshefte für Mathematik und Physik}}~22 (1911), pp. 137--49。}它也常被称为“与非”，因为它的特征真值表正是“$\eand$”的真值表的否定。
\factoidbox{\label{prop:upexpressive}“$\uparrow$” 自身就是功能完备的。}

功能完备性定理告诉我们，“$\enot$”和“$\eor$”是联合功能完备的。因此，只需证明：给定任何一个只包含这两个联结词的图式，我们都能将其重写为一个只包含“$\uparrow$”的等价图式。于是，与功能完备性定理辅助结果的证明类似，我们只需对辅助结果~\ref*{expressive:eor} 应用以下等价式：
		\begin{align*}
			\enot \metav{A} &\Leftrightarrow (\metav{A} \uparrow \metav{A})\\
			(\metav{A} \eor \metav{B}) & \Leftrightarrow ((\metav{A} \uparrow \metav{A}) \uparrow (\metav{B} \uparrow \metav{B}))
		\end{align*}

类似地，我们可以考虑联结词“$\downarrow$”：

\begin{center}
\begin{tabular}{c c | c}
$\metav{A}$ & $\metav{B}$ & $\metav{A} \mathrel{\downarrow} \metav{B}$\\
\hline
 T & T & F \\
 T & F & F  \\
 F & T & F  \\
 F & F & T
\end{tabular}
\end{center}

这有时被称为“皮尔斯箭头”（皮尔士本人称其为“ampheck”）。不过，更常见的名称是“或非”，因为它的特征真值表是“$\eor$”的否定，即表示“既不……也不……”。
	\factoidbox{
	“$\downarrow$” 自身就是功能完备的。}

与前面关于 $\uparrow$ 的结果证明类似，只不过这里应用的是以下等价式：
		\begin{align*}
			\enot \metav{A} &\Leftrightarrow (\metav{A} \downarrow \metav{A})\\
			(\metav{A} \eand \metav{B}) & \Leftrightarrow ((\metav{A} \downarrow \metav{A}) \downarrow (\metav{B} \downarrow \metav{B}))
		\end{align*}
以及辅助结果~\ref*{expressive:eand}。

\section{功能不完备的情况}

实际上，\emph{仅有的}两种自身功能完备的二元联结词是“$\uparrow$”和“$\downarrow$”。但我们如何证明这一点呢？更一般地说，我们如何证明某些联结词\emph{并非}联合功能完备？

显然，我们需要尝试找出某个真值表，它\emph{无法}仅用给定的联结词来表达。但这需要一些技巧。

为具体起见，让我们考虑“$\eor$”是否自身功能完备的问题。稍加思考便会清楚，并非如此。具体来说，任何仅包含析取的图式，其真值表都不可能和否定式相同，即：

\begin{center}
				\begin{tabular}{c | c}
				$\metav{A}$ & $\enot \metav{A}$\\
				\hline
				 T & F \\
				 F & T
				\end{tabular}
				\end{center}

其直观原因很简单：目标真值表的首行需要取值为假；但任何\emph{仅}包含$\eor$的图式的真值表，其首行总是为真。对于$\eand$、$\eif$和$\eiff$，情况同样如此。
 	\factoidbox{
		“$\eor$”、“$\eand$”、“$\eif$”和“$\eiff$”自身都不是功能完备的。}

事实上，以下结论成立：

\factoidbox{自身功能完备的\emph{仅有的}二元联结词是“$\uparrow$”和“$\downarrow$”。}

当然，这比针对原始联结词的证明要困难。例如，“排斥或”联结词的特征真值表首行并非T，因此上述方法不足以证明它不能表达所有真值表。要证明诸如“$\eiff$”和“$\enot$”\emph{联合}起来也非功能完备，同样更为困难。

\chapter{证明等值式}\label{ch:equivalences}

\section{等值式的可替换性}

回忆一下，从\cref{sec:equivalent}可知，$\metav{P}$ 和 $\metav{Q}$（在真值函项逻辑中）是等价的 \ifeff{} 对每个赋值而言，它们的真值都相同。我们已经见过很多这样的例子，并且用真值表和自然演绎证明来确立这种等价关系。在\cref{c:NormalForms}中，我们甚至证明了每个真值函项逻辑语句都等价于一个合取范式语句和一个析取范式语句。如果 $\metav{P}$ 和~$\metav{Q}$ 等价，那么它们总是具有相同的真值，任意一个都蕴涵另一个，并且从任意一个出发都可以证明另一个。

当然，等价的语句并非同一个语句：语句 $\enot\enot A$ 和~$A$ 可能总是有相同的真值，但第一个以 “$\enot$” 符号开头，而第二个则不然。不过你可能会想，是否总是能用一对等价语句中的一个替换另一个，并且得到的结果也是等价的呢？例如，考虑 $\enot\enot A \eif B$ 和 $A \eif B$。第二个语句正是将第一个语句中的 “$\enot\enot A$” 替换为~“$A$” 得到的。而这两个语句也确实等价。

这是一个普遍事实，并且不难理解原因。在任何赋值下，我们“从内到外”计算一个语句的真值。所以，当要决定 “$\enot\enot A \eif B$” 的真值时，我们首先计算 “$\enot\enot A$” 的真值，整个语句的真值随后就只取决于这个真值（可能是真或假）和语句的其余部分（即~“$B$” 的真值和 “$\eif$” 的真值表）。但是，由于 “$\enot\enot A$” 和~“$A$” 等价，在给定的赋值下它们总是具有相同的真值——因此，用~“$A$” 替换 “$\enot\enot A$” 不会改变整个语句的真值。当然，对于任何其他与 “$\enot\enot A$” 等价的语句，比如 “$A \eand (A \eor A)$”，情况也是如此。

为了一般性地陈述这个结果，我们使用记号 $\metav{R}(\metav{P})$ 来表示一个包含语句 $\metav{P}$ 作为其组成部分的语句。那么 $\metav{R}(\metav{Q})$ 就表示将 $\metav{P}$ 的出现替换为语句~$\metav{Q}$ 后得到的结果。例如，如果 $\metav{P}$ 是语句字母~“$A$”，$\metav{Q}$ 是语句~“$\enot\enot A$”，而 $\metav{R}(\metav{P})$ 是 “$A \eif B$”，那么 $\metav{R}(\metav{Q})$ 就是 “$\enot\enot A \eif B$”。

\factoidbox{如果 $\metav{P}$ 和 $\metav{Q}$ 是等价的，那么 $\metav{R}(\metav{P})$ 和 $\metav{R}(\metav{Q})$ 也是等价的。}

由此事实可知，任何形如 $\metav{R}(\metav{P}) \eiff \metav{R}(\metav{Q})$ 的语句（其中 $\metav{P}$ 和 $\metav{Q}$ 等价）都是重言式。然而，对于不同的~$\metav{R}$，其自然演绎证明将会有很大差异。（作为练习，请给出证明以表明 \begin{align*}
	& \vdash (\enot\enot P \eif Q) \eiff (P \eif Q)  \text{ 和}\\
	& \vdash (\enot\enot P \eand Q) \eiff (P \eand Q)
\end{align*} 是重言式，并比较这两个证明。）

还有另一个事实：如果两个语句 $\metav{P}$ 和~$\metav{Q}$ 是等价的，并且你在 $\metav{P}$ 和 $\metav{Q}$ 中用同一个语句~$\metav{R}$ 替换了某个相同的语句字母，那么得到的结果也是等价的。例如，如果你在 “$A \land B$” 和 “$B \land A$” 两者中都把 “$A$” 替换为，比方说，“$\enot C$”，你会得到 “$\enot C \land B$” 和 “$B \land \enot C$”，而这两者是等价的。我们也可以记录下这一点：

\factoidbox{等价关系在语句字母替换下保持不变，即，如果 $\metav{P}(A)$ 和 $\metav{Q}(A)$ 都包含语句字母~“$A$” 并且是等价的，那么语句 $\metav{P}(\metav{R})$ 和 $\metav{Q}(\metav{R})$（通过在两者中将 “$A$” 替换为 $\metav{R}$ 得到）也是等价的。}

这意味着，一旦我们证明了两条语句是等价的（例如，“$\enot\enot A$” 和 “$A$”，或 “$A \land B$” 和 “$B \land A$”），我们就知道它们所有共同的“实例”也都是等价的。请注意，我们从真值表或自然演绎证明中并不能直接得到这一点。例如，用于证明 “$\enot\enot A$” 和~“$A$” 等价的真值表，并\emph{不}能同时证明 “$\enot\enot(B \eif C)$” 和 “$B \eif C$” 是等价的：前者只需要 2 行，而后者需要~4 行。

\section{等价链}

当你想验证两个语句等价时，当然可以做真值表，或者寻找形式证明。但还有一种更简单的方法，它基于我们刚刚讨论的等价替换原则：依据一个由简单等价关系构成的小型目录，将第一个语句中的部分替换为等价的部分，并不断重复此过程，直到得到第二个语句。

这种证明语句等价的方法是由前一节的两个事实所支撑的。第一个事实告诉我们，\emph{如果}，比如说，$\enot\enot \metav{P}$ 和 $\metav{P}$ 是等价的（对于任意语句 $\metav{P}$），那么在一个语句中用 $\metav{P}$ 替换 $\enot\enot\metav{P}$ 会产生一个等价的语句。第二个事实告诉我们，对于任意语句 $\metav{P}$，$\enot\enot \metav{P}$ 和 $\metav{P}$ \emph{总是}等价的。（一个简单的真值表就能表明 “$\enot\enot A$” 和 “$A$” 是等价的。）根据第二个事实我们知道，无论何时我们将 “$\enot\enot A$” 和 “$A$” 中的 “$A$” 替换为某个语句 $\metav{P}$，得到的两个语句都是等价的。换句话说，从 “$\enot\enot A$” 和 “$A$” 等价以及第二个事实，我们可以得出结论：对于任意语句 $\metav{P}$，$\enot\enot \metav{P}$ 和 $\metav{P}$ 是等价的。

我们来看一个例子。根据德摩根律，下列成对的语句是等价的：
\begin{align*}
	\enot (A \eand B) & \Leftrightarrow (\enot A \eor \enot B)\\
	\enot (A \eor B) & \Leftrightarrow (\enot A \eand \enot B)
\end{align*}
这可以通过构造两个真值表来验证，或通过四个自然演绎证明来表明：
\begin{align*}
	\enot (A \eand B) & \vdash (\enot A \eor \enot B)\\
	(\enot A \eor \enot B) & \vdash \enot (A \eand B)\\
	\enot (A \eor B) & \vdash (\enot A \eand \enot B)\\
	(\enot A \eand \enot B) & \vdash \enot (A \eor B)
\end{align*}

根据第二个事实，\emph{任何}具有下列形式的语句对都是等价的：
\begin{align*}
	\enot (\metav{P} \eand \metav{Q}) & \Leftrightarrow (\enot \metav{P} \eor \enot \metav{Q})\\
	\enot (\metav{P} \eor \metav{Q}) & \Leftrightarrow (\enot \metav{P} \eand \enot \metav{Q})
\end{align*}
现在考虑语句 “$\enot(R \eor (S \eand T))$”。我们将找到一个等价语句，其中所有的 “$\enot$” 符号都直接作用于语句字母。第一步，我们把它看作形式为$\enot(\metav{P} \eor \metav{Q})$的语句——那么 $\metav{P}$ 就是语句 “$R$”，而 $\metav{Q}$ 是 “$(S \eand T)$”。由于$\enot(\metav{P} \eor \metav{Q})$等价于$(\enot \metav{P} \eand \enot \metav{Q})$（根据德摩根律的第二条），我们可以将整个语句替换为$(\enot \metav{P} \eand \enot \metav{Q})$。在这个例子中（其中 $\metav{P}$ 是 “$R$”，$\metav{Q}$ 是 “$(S \eand T)$”）我们得到 “$(\enot R \eand \enot (S \eand T))$”。这个新语句包含 “$\enot (S \eand T)$” 作为一个部分。它属于$\enot(\metav{P} \eand \metav{Q})$这种形式，只不过现在 $\metav{P}$ 是语句字母 “$S$”，$\metav{Q}$ 是 “$T$”。根据德摩根律（这次是第一条），它等价于$(\enot\metav{P} \eor \enot\metav{Q})$，或者在这个具体例子中，等价于 “$(\enot S \eor \enot T)$”。于是我们可以用 “$(\enot S \eor \enot T)$” 替换 “$\enot (S \eand T)$” 这个部分。这样就得到了语句 “$(\enot R \eand (\enot S \eor \enot T))$”，其中所有的 “$\enot$” 符号都直接作用于语句字母。我们已经尽可能地将否定符号“向内推”了。我们可以通过列出各个步骤来记录这样一条等价链，就像在自然演绎中一样，记录每一步使用了哪个基本等价关系：
\begin{align*}
	& \fbox{\enot(\fbox{$R$} \eor \fbox{$(S \eand T)$})} \\
	& \fbox{(\enot(\fbox{$R$} \eand \enot\fbox{$(S \eand T)$})} && \text{DeM}\\
	& (\enot(R \eand \fbox{$\enot(\fbox{$S$} \eand \fbox{$T$})$}) \\
	& (\enot(R \eand \fbox{$(\enot \fbox{$S$} \eor \enot \fbox{$T$})$}) && \text{DeM}
\end{align*}
为了清晰起见，我们标出了被替换的语句以及那些与德摩根律中的 $\metav{P}$ 和~$\metav{Q}$ 相匹配的部分，但这并非必要，后文也不会一直这样做。

在\cref{tab:equivalences}中，我们给出了一份基本等价关系列表，可用于构建这样的等价链。其标签缩写自相应逻辑定律的习惯名称：双重否定（DN）、德摩根（DeM）、交换律（Comm）、分配律（Dist）、结合律（Assoc）、幂等律（Id）和吸收律（Abs）。

\begin{table}
	\begin{align*}
\lnot\lnot \metav{P} & \Leftrightarrow \metav{P} && \text{(DN)}\\[2ex]
(\metav{P} \eif \metav{Q}) & \Leftrightarrow (\lnot \metav{P} \lor \metav{Q})
&& \text{(Cond)}\\
\lnot(\metav{P} \eif \metav{Q}) & \Leftrightarrow (\metav{P} \land \lnot\metav{Q}) \\[2ex]
(\metav{P} \eiff \metav{Q}) & \Leftrightarrow ((\metav{P} \eif \metav{Q}) \land  (\metav{Q} \eif \metav{P}))
&& \text{(Bicond)}\\[2ex]
\lnot(\metav{P} \land \metav{Q}) & \Leftrightarrow (\lnot\metav{P} \lor \lnot\metav{Q})
&& \text{(DeM)}\\
\lnot(\metav{P} \lor \metav{Q}) & \Leftrightarrow (\lnot\metav{P} \land \lnot\metav{Q}) \\[2ex]
(\metav{P} \lor \metav{Q}) & \Leftrightarrow (\metav{Q} \lor \metav{P}) & &\text{(Comm)}\\
(\metav{P} \land \metav{Q}) & \Leftrightarrow (\metav{Q} \land \metav{P})\\[2ex]
(\metav{P} \land (\metav{Q} \lor \metav{R})) & \Leftrightarrow ((\metav{P} \land \metav{Q}) \lor (\metav{P} \land \metav{R}))
&& \text{(Dist)}\\
(\metav{P} \lor (\metav{Q} \land \metav{R})) & \Leftrightarrow ((\metav{P} \lor \metav{Q}) \land (\metav{P} \lor \metav{R}))\\
(\metav{P} \lor (\metav{Q} \lor \metav{R})) & \Leftrightarrow ((\metav{P} \lor \metav{Q}) \lor \metav{R}) & &\text{(Assoc)}\\
(\metav{P} \land (\metav{Q} \land \metav{R})) & \Leftrightarrow ((\metav{P} \land \metav{Q}) \land \metav{R})\\[2ex]
(\metav{P} \lor \metav{P}) & \Leftrightarrow \metav{P} && \text{(Id)}\\
(\metav{P} \land \metav{P}) & \Leftrightarrow \metav{P}\\[2ex]
(\metav{P} \land (\metav{P} \lor \metav{Q})) & \Leftrightarrow \metav{P}&& \text{(Abs)}\\
(\metav{P} \lor (\metav{P} \land \metav{Q})) & \Leftrightarrow \metav{P}\\[2ex]
(\metav{P} \land (\metav{Q} \lor \lnot\metav{Q})) & \Leftrightarrow \metav{P}  && \text{(Simp)}\\
(\metav{P} \lor (\metav{Q} \land \lnot\metav{Q})) & \Leftrightarrow \metav{P}\\[2ex]
(\metav{P} \lor (\metav{Q} \lor \lnot\metav{Q})) & \Leftrightarrow (\metav{Q} \lor \lnot\metav{Q}) \\
(\metav{P} \land (\metav{Q} \land \lnot\metav{Q})) & \Leftrightarrow (\metav{Q} \land \lnot\metav{Q})
\end{align*}
\caption{基本等价关系}
\label{tab:equivalences}
\end{table}

\section{寻找等价范式}

在\cref{c:NormalForms}中我们证明了，真值函项逻辑的每个语句都等价于一个析取范式语句和一个合取范式语句。我们当时提供了一种方法：先构造真值表，然后从真值表中“读出”一个等价于原语句的析取范式或合取范式语句。这种方法有两个缺点。第一个是得到的析取范式或合取范式语句并不总是最短的。第二个是当初始语句包含超过少数几个语句字母时，该方法本身就变得难以应用（因为包含 $n$ 个语句字母的语句，其真值表有 $2^n$ 行）。

我们可以将等价链用作另一种方法：为了找到析取范式语句，我们可以连续应用基本等价关系，直到找到一个等价的、且是析取范式的语句。回想一下析取范式语句必须满足的条件：

\begin{compactlist}
	\item[(\textsc{dnf1})] 语句中出现的联结词仅限于否定、合取和析取；
	\item[(\textsc{dnf2})] 每个否定符号的辖域都是极小的（即，任何 “$\enot$” 后面都紧跟着一个原子语句）；
	\item[(\textsc{dnf3})] 任何合取的辖域内部都不出现析取。
\end{compactlist}

条件 (\textsc{dnf1}) 要求我们从语句中消除所有 “$\eif$” 和～“$\eiff$” 符号。这就要用到基本等价规则 (Cond) 和 (Bicond)。例如，假设我们从语句
\begin{align*}
	& \enot(A \eand \enot C) \land (\enot A\eif \enot B).
\intertext{开始。我们可以使用 (Cond) 规则消去 “$\eif$”。这里 $\metav{P}$ 是 “$\enot A$”，$\metav{Q}$ 是 “$\enot{B}$”。于是得到：}
& \enot(A \eand \enot C) \land (\enot\enot A \eor \enot B) && \text{Cond}
\intertext{双重否定可以消去，因为 “$\enot\enot A$” 等价于～“$A$”：}
& \enot(A \eand \enot C) \land (A \eor \enot B) && \text{DN}
\intertext{现在条件 (\textsc{dnf1}) 已满足：我们的语句只包含 “$\enot$”、“$\eand$” 和 “$\eor$”。条件 (\textsc{dnf2}) 要求所有 “$\enot$” 都直接作用于语句字母。但第一个合取支尚未满足此条件。为确保满足 (\textsc{dnf2})，我们根据需要多次使用德摩根律和双重否定 (DN) 律。}
& (\enot A \eor \enot\enot C) \land (A \eor \enot B) && \text{DeM}\\
& (\enot A \eor C) \land (A \eor \enot B) && \text{DN}
\intertext{得到的语句现在处于合取范式——它是语句字母及其否定的析取的合取。但我们需要的是析取范式，即满足 (\textsc{dnf3}) 的语句：没有 “$\eor$” 出现在 “$\eand$” 的辖域内。我们使用分配律 (Dist) 来达成这一点。最后一个语句形如 $\metav{P} \land (\metav{Q} \lor \metav{R})$，其中 $\metav{P}$ 是 “$(\enot A \eor C)$”，$\metav{Q}$ 是 “$A$”，$\metav{R}$ 是 “$\enot B$”。应用一次 (Dist) 得到：}
& ((\enot A \eor C) \eand A)) \eor ((\enot A \eor C) \eand \enot B) && \text{Dist}
\intertext{虽然看起来更复杂了，但只要继续做下去就会好转！这两个析取支看起来几乎可以再次应用 (Dist)，只是 “$\eor$” 所在的位置不对。这时就需要用到交换律 (Comm)。我们将其应用于 “$(\enot A \eor C) \eand A$”：}
& (A \eand (\enot A \eor C)) \eor ((\enot A \eor C) \eand \enot B) && \text{Comm}
\intertext{我们可以对结果部分 “$A \eand (\enot A \eor C)$” 再次应用 (Dist)：}
& ((A \eand \enot A) \eor (A \eand C)) \eor ((\enot A \eor C) \eand \enot B) && \text{Dist}
\intertext{现在左半部分没有 “$\eor$” 出现在 “$\eand$” 的辖域内。让我们对右半部分应用相同的原理：}
& ((A \eand \enot A) \eor (A \eand C)) \eor (\enot B \eand (\enot A \eor C)) && \text{Comm}\\
& ((A \eand \enot A) \eor (A \eand C)) \eor ((\enot B \eand \enot A) \eor (\enot B \eand C)) && \text{Dist}
\intertext{我们的语句现在已经是析取范式了！但我们还可以稍作简化：“$(A \eand \enot A)$” 在真值函项逻辑中是一个矛盾式，即它总是假的。而将一个总是假的公式通过 “$\eor$” 与任意语句 $\metav{P}$ 结合，所得结果等价于 $\metav{P}$ 本身。这正是简化规则 (Simp) 的第二条。}
& ((A \eand C) \eor (A \eand \enot A)) \eor ((\enot B \eand \enot A) \eor (\enot B \eand C)) && \text{Comm}\\
& (A \eand C) \eor ((\enot B \eand \enot A) \eor (\enot B \eand C)) && \text{Simp}
\end{align*}
最终结果仍是析取范式，并且更简洁了一些。它也远比我们通过 \cref{c:NormalForms} 中方法得到的析取范式简单。事实上，我们起始的语句可能就是 \cref{s:DNFTruthTable} 中的 $\metav{S}$——它恰好具有那里用作示例的真值表。我们在那里找到的析取范式（\ifHTMLtarget 在 ~\cref{s:DNFTruthTable} 节\else 见第 ~\pageref{longDNF} 页\fi）是（带有所有必需的括号）：
$$((((A \eand B) \eand C) \eor ((A \eand \enot B) \eand C)) \eor ((\enot A \eand \enot B) \eand C)) \eor ((\enot A \eand \enot B) \eand \enot C)$$

\practiceproblems
\problempart
\label{pr.DNF2}
考虑以下语句：

\begin{compactlist}
	\item $(A \eif \enot B)$
	\item $\enot (A \eiff B)$
	\item $(\enot A \eor \enot (A \eand B))$
	\item $(\enot (A \eif B ) \eand (A \eif C))$
	\item $(\enot (A \eor B) \eiff ((\enot C \eand \enot A) \eif \enot B))$
	\item $((\enot (A \eand \enot B) \eif C) \eand \enot (A \eand D))$
\end{compactlist}

对每个语句，通过给出一系列等价变换，分别找到一个等价的析取范式语句和一个等价的合取范式语句。使用 (Id)、(Abs) 和 (Simp) 规则尽可能简化你的语句。

\chapter{可靠性}\label{ch:Soundness}

在本章中，我们将真值函项逻辑的语义与其自然演绎 \emph{证明系统}（如 \cref{ch.NDTFL} 所定义）联系起来。我们将证明这个形式证明系统是安全的：你只能从一组前提 $\Gamma$ 出发，证明那些真正被 $\Gamma$ 所蕴涵的结论。
直观上，一个形式证明系统是可靠的 \ifeff{} 它不允许你证明任何无效论证。这显然是一个极其理想的性质。它告诉我们，我们的证明系统永远不会将我们引入歧途。事实上，如果我们的证明系统不可靠，那么我们就无法信任我们的任何证明。本章的目标正是证明我们的证明系统是可靠的。

我们来更精确地阐述这个想法。我们将使用希腊字母 $\Gamma$（gamma）来缩写一个语句的集合。一个形式证明系统是 \define{可靠的}（相对于给定的语义）\emph{\ifeff}，只要存在以 $\Gamma$ 中语句为（未解除的）假设而推导出 $\metav{C}$ 的形式证明，那么 $\Gamma$ 就真正蕴涵 $\metav{C}$（根据该语义）。换言之，要证明真值函项逻辑的证明系统是可靠的，我们需要证明下述定理：

\begin{factoidboxe}\textbf{可靠性定理。} 对于任意语句集 $\Gamma$ 和语句 $\metav{C}$：若 $\Gamma\proves\metav{C}$，则 $\Gamma \entails\metav{C}$。
\end{factoidboxe}

为证明此定理，我们将逐一检视真值函项逻辑证明系统中的每条规则。我们希望表明，这些规则的任何应用都不会将我们引入歧途。既然一个证明仅仅是这些规则的重复应用，那么任何证明也就都不会将我们引入歧途。至少，这是大体的思路。

首先，我们必须更精确地界定何谓“引入歧途”。我们称证明中的某一行是 \define{闪亮的} \ifeff{} 该行所依赖的假设蕴涵了该行上的语句。\footnote{“闪亮”一词并非逻辑学家的标准术语。} 为说明此概念，考虑如下证明：


\begin{fitchproof}
		\hypo{fgh}{F\eif(G\eand H)}\PR
		\open
			\hypo{f}{F}\AS
			\have{gh}{G \eand H}\ce{fgh,f}
			\have{g}{G}\ae{gh}
		\close
		\have{fg}{F \eif G}\ci{f-g}
	\end{fitchproof}

\noindent\noindent
第 $1$ 行是闪亮的 \ifeff{} $F \eif (G \eand H) \entails F \eif (G \eand H)$。你应当能轻易信服，第 $1$ 行的确是闪亮的！类似地，第 $4$ 行是闪亮的 \ifeff{} $F \eif (G \eand H), F \entails G$。同样容易验证第 $4$ 行是闪亮的。事实上，该真值函项逻辑证明中的每一行都是如此。我们希望证明这并非巧合。也就是说，我们希望证明：


\begin{factoidboxe}\textbf{闪亮引理。}
		每个真值函项逻辑证明中的每一行都是闪亮的。
	\end{factoidboxe}

\noindent
由此我们将知道，在证明的任何一行我们都未曾走入歧途。事实上，基于闪亮引理，可靠性定理的证明将很简单：

\emph{证明。} 假设 $\Gamma \proves \metav{C}$。则存在一个真值函项逻辑证明，其中 $\metav{C}$ 出现在最后一行，且其所有未解除的假设均属于 $\Gamma$。闪亮引理断言每个真值函项逻辑证明的每一行都是闪亮的。故这最后一行是闪亮的，即 $\Gamma \entails \metav{C}$。证毕。

现在剩下的任务是证明闪亮引理。

为此，我们注意到任何真值函项逻辑证明中的每一行，要么是前提或假设，要么是通过应用某条规则得到的。前提天然是闪亮的：若 $\metav{A}$ 是前提，则它属于 $\Gamma$ 中的语句，显然有 $\Gamma \entails \metav{A}$。假设也是闪亮的，因为假设 $\metav{A}$ 仅依赖于其自身，而 $\metav{A}$ 显然蕴涵 $\metav{A}$。所以我们需要证明的是，应用真值函项逻辑证明系统中的任何规则都不会导致我们出错。更精确地说，我们称一条推理规则是 \define{规则可靠的} \emph{\ifeff} 对于所有真值函项逻辑证明，如果我们通过应用该规则在证明中得到一行，并且该证明中所有先前的行都是闪亮的，那么我们的新行也是闪亮的。我们需要证明，真值函项逻辑证明系统中的\emph{每一条}规则都是规则可靠的。

我们将在下文完成此事。但一旦证明了每条规则都是规则可靠的，闪亮引理便立刻可得：

\emph{证明。} 任取一个真值函项逻辑证明，从第 $1$ 行开始。它必是前提或假设，而如上所述，这些都是闪亮的。考虑第 $2$ 行。若它是前提或假设，则它是闪亮的。否则，它是由第 $1$ 行应用一条规则可靠的推理规则得到的。既然第 $1$ 行闪亮，第 $2$ 行也闪亮。接着考虑第 $3$ 行。若它是前提或假设，则它是闪亮的。否则，它是由某个（或某些）先前行应用一条规则可靠的推理规则得到的，而我们已经确定所有先前行都是闪亮的。因此，第 $3$ 行也是闪亮的。
依此类推。一般地，第 $n$ 行上的语句要么是前提或假设（因而是闪亮的），要么是通过应用一条规则可靠的形式推理规则从某些闪亮的行得到的。由于所有先前的行都是闪亮的，故第 $n$ 行本身也是闪亮的。对任意真值函项逻辑证明，我们均可按此推理，从第 $1$ 行迭代至最后一行，从而证明该证明的每一行都是闪亮的。证毕。

剩下的便是证明每条规则都是规则可靠的。这并不困难，但相当耗时，因为我们需要逐一检查每条规则，而真值函项逻辑的证明系统包含大量规则！为略微加快进程，我们引入一个方便的缩写：用 “$\Delta_i$” (“delta”) 表示在当前考虑的真值函项逻辑证明中，第 $i$ 行所依赖的那些假设的集合（上下文会指明是哪个证明）。这包括该证明的所有前提，以及在第 $i$ 行时仍未解除的所有子证明假设。我们首先将上述关于前提和假设的观察记录如下。

\begin{factoidboxe}
	真值函项逻辑证明中的前提和假设是闪亮的。
\end{factoidboxe}

如果 $\metav{A}$ 是第 $n$ 行的前提，那么它属于 $\Delta_n$，因为 $\Delta_n$ 包含了该证明的所有前提。如果它是在第 $n$ 行作为子证明的假设引入的，那么子证明中的所有内容（包括第 $n$ 行，即 $\metav{A}$ 本身）都依赖于 $\metav{A}$，因此 $\metav{A}$ 属于 $\Delta_n$。无论哪种情况，都有 $\Delta_n \entails \metav{A}$。

接下来，我们证明所有推理规则都是规则可靠的。

\begin{factoidboxe}
$\eand$I 是规则可靠的。
\end{factoidboxe}

\emph{证明.} 考虑任意真值函项逻辑证明中任意一次 $\eand$I 的应用，即类似于以下的形式：
\begin{fitchproof}
	\have[i]{a}{\metav{A}}
	\have[j]{b}{\metav{B}}
	\have[n]{c}{\metav{A}\eand\metav{B}} \ai{a, b}
\end{fitchproof}
\noindent
为了证明 $\eand$I 是规则可靠的，我们假设第 $n$ 行之前的每一行都是闪亮的；我们的目标是证明第 $n$ 行也是闪亮的，即 $\Delta_n \entails \metav{A} \eand \metav{B}$。

因此，设 $v$ 为任意使得 $\Delta_{n}$ 中所有语句都为真的赋值。

我们首先证明 $v$ 使 $\metav{A}$ 为真。为此，注意 $\Delta_i$ 中的所有语句都属于 $\Delta_{n}$。根据假设，第 $i$ 行是闪亮的。所以，任何使得 $\Delta_i$ 中所有语句为真的赋值都使 $\metav{A}$ 为真。由于 $v$ 确实使得 $\Delta_i$ 中所有语句为真，因此 $v$ 也使 $\metav{A}$ 为真。

类似地，我们可以知道 $v$ 使 $\metav{B}$ 为真。

因此，$v$ 使 $\metav{A}$ 为真且 $v$ 使 $\metav{B}$ 为真。于是，$v$ 使 $\metav{A}\eand\metav{B}$ 为真。所以，任何使得 $\Delta_{n}$ 中所有语句为真的赋值也都使 $\metav{A} \eand \metav{B}$ 为真。也就是说：第 $n$ 行是闪亮的。\textsc{QED}

接下来建立规则可靠性的所有引理，实质上都将具有与此相同的基本结构。

\begin{factoidboxe}
$\eand$E 是规则可靠的。
\end{factoidboxe}

\emph{证明.}
假设在某个真值函项逻辑证明中，第 $n$ 行之前的每一行都是闪亮的，并且第 $n$ 行使用了 $\eand$E。那么情况如下：
\begin{fitchproof}
			   \have[i]{ab}{\metav{A}\eand\metav{B}}
			   \have[n]{a}{\metav{A}} \ae{ab}
		   \end{fitchproof}
\noindent
（也可能在第 $n$ 行得到的是 $\metav{B}$；但类似的推理同样适用）。设 $v$ 为任意使得 $\Delta_{n}$ 中所有语句都为真的赋值。注意 $\Delta_i$ 中的所有语句都属于 $\Delta_{n}$。根据假设，第 $i$ 行是闪亮的。所以，任何使得 $\Delta_i$ 中所有语句为真的赋值都使 $\metav{A}\eand\metav{B}$ 为真。因此，$v$ 使 $\metav{A}\eand\metav{B}$ 为真，从而 $v$ 使 $\metav{A}$ 为真。所以 $\Delta_{n} \entails \metav{A}$。\textsc{QED}

\begin{factoidboxe}
$\eor$I 是规则可靠的。
\end{factoidboxe}

我们将其留作练习。

\begin{factoidboxe}
$\eor$E 是规则可靠的。
\end{factoidboxe}

\emph{证明.}
假设在某个真值函项逻辑证明中，第 $n$ 行之前的每一行都是闪亮的，并且第 $n$ 行使用了 $\eor$E。那么情况如下：
\begin{fitchproof}
	   \have[m]{aob}{\metav{A}\eor\metav{B}}
	   \open
		   \hypo[i]{a}{\metav{A}}\AS %\by{want \metav{C}}{}
		   \have[j]{c1}{\metav{C}}
	   \close
	   \open
		   \hypo[k]{b}{\metav{B}}\AS %\by{want \metav{C}}{}
		   \have[l]{c2}{\metav{C}}
	   \close
	   \have[n]{ab}{\metav{C}}\oe{aob, a-c1,b-c2}
   \end{fitchproof}
\noindent
设 $v$ 为任意使得 $\Delta_{n}$ 中所有语句都为真的赋值。注意 $\Delta_m$ 中的所有语句都属于 $\Delta_{n}$。根据假设，第 $m$ 行是闪亮的。所以任何使得 $\Delta_{n}$ 中所有语句为真的赋值都使 $\metav{A} \eor \metav{B}$ 为真。特别地，$v$ 使 $\metav{A} \eor \metav{B}$ 为真，因此要么 $v$ 使 $\metav{A}$ 为真，要么 $v$ 使 $\metav{B}$ 为真。我们现在分两种情况讨论：
\begin{numberlist}
	   \item \emph{$v$ 使 $\metav{A}$ 为真。} $\Delta_i$ 中的所有语句都属于 $\Delta_{n}$，但可能不包含 $\metav{A}$。由于 $v$ 使得 $\Delta_{n}$ 中所有语句为真，并且也使 $\metav{A}$ 为真，所以 $v$ 使得 $\Delta_i$ 中所有语句为真。现在，根据假设，第 $j$ 行是闪亮的；所以 $\Delta_{j} \entails \metav{C}$。但 $\Delta_i$ 正好就是 $\Delta_{j}$，因此 $\Delta_i \entails \metav{C}$。所以，任何使得 $\Delta_i$ 中所有语句为真的赋值都使 $\metav{C}$ 为真。而 $v$ 正是这样一个赋值。因此 $v$ 使 $\metav{C}$ 为真。
	   \item \emph{$v$ 使 $\metav{B}$ 为真。} 用完全相同的方式推理，考虑第 $k$ 行和第 $l$ 行，$v$ 使 $\metav{C}$ 为真。
\end{numberlist}
无论哪种情况，$v$ 都使 $\metav{C}$ 为真。所以 $\Delta_n \entails \metav{C}$。
\textsc{QED}

\begin{factoidboxe}
	$\enot$E 是规则可靠的。
\end{factoidboxe}

\emph{证明.}
假设在某个真值函项逻辑证明中，第 $n$ 行之前的每一行都是闪亮的，并且第 $n$ 行使用了 $\enot$E。那么情况如下：
\begin{fitchproof}
   \have[i]{i}{\metav{A}}
   \have[j]{j}{\enot\metav{A}}
   \have[n]{nb}{\ered}\ri{i, j}
\end{fitchproof}
\noindent
注意：所有 $\Delta_i$ 和所有 $\Delta_j$ 都在 $\Delta_{n}$ 之中。根据假设，第 $i$ 行和第 $j$ 行都是闪亮的。因此，任何使 $\Delta_{n}$ 中所有语句为真的赋值都必定会使 $\metav{A}$ 和 $\enot\metav{A}$ 同时为真。但没有任何赋值能做到这一点。所以，没有赋值能使 $\Delta_{n}$ 中所有语句为真。于是，$\Delta_{n} \entails \ered$（空乏地成立）。\textsc{QED}

\begin{factoidboxe}
	X 是规则可靠的。
\end{factoidboxe}

我们将其留作练习。

\begin{factoidboxe}
	$\enot$I 是规则可靠的。
\end{factoidboxe}

\emph{证明.}
假设在某个真值函项逻辑证明中，第 $n$ 行之前的每一行都是闪亮的，并且第 $n$ 行使用了 $\enot$I。那么情况如下：
\begin{fitchproof}
   \open
	   \hypo[i]{a}{\metav{A}}\AS
	   \have[j]{b}{\ered}
   \close
   \have[n]{na}{\enot\metav{A}}\ni{a-b}
\end{fitchproof}
\noindent
令 $v$ 为任意使 $\Delta_{n}$ 中所有语句为真的赋值。注意，所有 $\Delta_{n}$ 中的语句都在 $\Delta_i$ 之中，可能的例外是 $\metav{A}$ 本身。根据假设，第 $j$ 行是闪亮的。但没有任何赋值能使 “$\ered$” 为真，因此也没有任何赋值能使 $\Delta_{j}$ 中所有语句为真。由于语句集 $\Delta_i$ 就是语句集 $\Delta_{j}$，所以没有赋值能使 $\Delta_i$ 中所有语句为真。既然 $v$ 使 $\Delta_{n}$ 中所有语句为真，那么它必然使 $\metav{A}$ 为假，从而使 $\enot \metav{A}$ 为真。因此，$\Delta_n \entails \enot \metav{A}$。\textsc{QED}

\begin{factoidboxe}\label{lem:LastRuleSound} IP、$\eif$I、$\eif$E、$\eiff$I 和 $\eiff$E 都是规则可靠的。
\end{factoidboxe}

我们将这些留作练习。

至此，我们已确立证明系统中所有基本规则都是规则可靠的。最后，我们证明：

\begin{factoidboxe}
证明系统中的所有派生规则都是规则可靠的。
\end{factoidboxe}

\emph{证明.}
假设我们在某个真值函项逻辑证明的第 $n$ 行使用了一个派生规则得到了某个语句 $\metav{A}$，并且所有先前的行都是闪亮的。任何对派生规则的使用都可以（以冗长为代价）用多次使用基本规则来替代。也就是说，我们本可以使用基本规则在某个行 $n + k$ 上写下 $\metav{A}$，而不引入任何额外的假设。因此，通过多次应用我们已有的结果——所有基本规则都是规则可靠的（实际上是 $k + 1$ 次）——我们可以看出第 $n+k$ 行是闪亮的。因此，该派生规则是规则可靠的。\textsc{QED}

就是这样！我们已经证明了每条规则——无论是基本的还是派生的——都是规则可靠的，而这正是我们建立闪亮引理，进而建立可靠性定理所需要的全部。

不过，如果在本章末尾再次非正式地解释一下我们所做的事情，或许会有所帮助。一个形式证明只是一系列——任意长度的——规则应用。我们已经表明，任何规则的任何应用都不会误导你。由此可见，没有任何形式证明会误导你。也就是说：我们的证明系统是可靠的。

\practiceproblems

\problempart
\label{pr.Soundness}
完成本章留作练习的引理。也就是说，证明以下规则是规则可靠的：
\begin{compactlist}
		\item $\eor$I。（\emph{提示}：与 $\eand$E 的情况类似。）
		\item X。（\emph{提示}：与 $\enot$E 的情况类似。）
		\item $\eif$I。（\emph{提示}：与 $\eor$E 的情况类似。）
		\item $\eif$E。
		\item IP。（\emph{提示}：与 $\enot$I 的情况类似。）
	\end{compactlist}
